\QHeading{画图解决问题。}
\par\noindent


\Problem{① 班主任调查班级学生周末活动情况如下：有36人去公园游玩，有30人去图书馆看书，有12人既去了公园又去了图书馆，还有5人两个地方都没去。这个班级一共有多少人？}
\par




\MethodTitle{方法一}
\Answer{{\Formula{\begin{aligned}[t]&\phantom{={}}36+30-12+5\\&=66-12+5\\&=54+5\\&=59\text{（人）}\end{aligned}}}\par 答：这个班一共有59人。}

\MethodTitle{方法二：分成四个互不重叠的部分}
\Answer{只去公园、两处都去、只去图书馆、两处都没去，四部分合起来。}
\Answer{\Formula{\begin{aligned}[t]&\phantom{={}}(36-12)+12+(30-12)+5\\
&=24+12+18+5\\
&=36+18+5\\
&=54+5\\
&=59\text{（人）}\end{aligned}}}
\Answer{答：这个班级一共有59人。}
\Problem{② 学校音乐兴趣小组成员至少要会演奏吉他和葫芦丝中的一种，其中会弹吉他的有25人，会吹葫芦丝的有19人，两种乐器都会的有5人。学校音乐兴趣小组一共有多少人？}



\MethodTitle{方法一}
\Answer{{\Formula{\begin{aligned}[t]&\phantom{={}}25+19-5\\&=44-5\\&=39\text{（人）}\end{aligned}}}\par 答：音乐兴趣小组一共有39人。}

\MethodTitle{方法二：先求只属于第二类的人数}
\Answer{第一类人数已经包含两类兼有的人，只需再加上仅属于第二类的人。}
\Answer{\Formula{\begin{aligned}[t]&\phantom{={}}25+(19-5)\\
&=25+14\\
&=39\text{（人）}\end{aligned}}}
\Answer{答：音乐兴趣小组一共有39人。}

\MethodTitle{方法三：分成互不重叠的三类}
\Answer{分别计数“只属于第一类”“两类兼有”“只属于第二类”，每人只数一次。}
\Answer{\Formula{\begin{aligned}[t]&\phantom{={}}(25-5)+5+(19-5)\\
&=20+5+14\\
&=25+14\\
&=39\text{（人）}\end{aligned}}}
\Answer{答：音乐兴趣小组一共有39人。}
\Problem{③ 四（1）班有42人，其中会骑自行车的有17人，会游泳的有14人，既会骑自行车又会游泳的有4人。两种运动都不会的有多少人？}



\MethodTitle{方法一}
\Answer{{\Formula{\begin{aligned}[t]&\phantom{={}}42-(17+14-4)\\&=42-(31-4)\\&=42-27\\&=15\text{（人）}\end{aligned}}}\par 答：两种运动都不会的有15人。}

\MethodTitle{方法二：先扣第一项，再扣只参加第二项的人}
\Answer{第一项已包括两项都参加的人，所以第二次只能扣除第二项独有的人。}
\Answer{\Formula{\begin{aligned}[t]&\phantom{={}}42-17-(14-4)\\
&=42-17-10\\
&=25-10\\
&=15\text{（人）}\end{aligned}}}
\Answer{答：两种运动都不会的有15人。}

\MethodTitle{方法三：先连减，再补回重复扣除的人}
\Answer{如果两项人数都减去，交叉部分被减了两次，需要加回来一次。}
\Answer{\Formula{\begin{aligned}[t]&\phantom{={}}42-17-14+4\\
&=25-14+4\\
&=11+4\\
&=15\text{（人）}\end{aligned}}}
\Answer{答：两种运动都不会的有15人。}
\Problem{④ 某班有40人，对数学感兴趣的有17人，对音乐感兴趣的有13人，对数学和音乐都感兴趣的有8人。对数学和音乐都不感兴趣的有多少人？}



\MethodTitle{方法一}
\Answer{{\Formula{\begin{aligned}[t]&\phantom{={}}40-(17+13-8)\\&=40-(30-8)\\&=40-22\\&=18\text{（人）}\end{aligned}}}\par 答：对数学和音乐都不感兴趣的有18人。}

\MethodTitle{方法二：先扣第一项，再扣只参加第二项的人}
\Answer{第一项已包括两项都参加的人，所以第二次只能扣除第二项独有的人。}
\Answer{\Formula{\begin{aligned}[t]&\phantom{={}}40-17-(13-8)\\
&=40-17-5\\
&=23-5\\
&=18\text{（人）}\end{aligned}}}
\Answer{答：对数学和音乐都不感兴趣的有18人。}

\MethodTitle{方法三：先连减，再补回重复扣除的人}
\Answer{如果两项人数都减去，交叉部分被减了两次，需要加回来一次。}
\Answer{\Formula{\begin{aligned}[t]&\phantom{={}}40-17-13+8\\
&=23-13+8\\
&=10+8\\
&=18\text{（人）}\end{aligned}}}
\Answer{答：对数学和音乐都不感兴趣的有18人。}
